% Options for packages loaded elsewhere
\PassOptionsToPackage{unicode}{hyperref}
\PassOptionsToPackage{hyphens}{url}
\documentclass[
]{article}
\usepackage{xcolor}
\usepackage{amsmath,amssymb}
\setcounter{secnumdepth}{-\maxdimen} % remove section numbering
\usepackage{iftex}
\ifPDFTeX
  \usepackage[T1]{fontenc}
  \usepackage[utf8]{inputenc}
  \usepackage{textcomp} % provide euro and other symbols
\else % if luatex or xetex
  \usepackage{unicode-math} % this also loads fontspec
  \defaultfontfeatures{Scale=MatchLowercase}
  \defaultfontfeatures[\rmfamily]{Ligatures=TeX,Scale=1}
\fi
\usepackage{lmodern}
\ifPDFTeX\else
  % xetex/luatex font selection
\fi
% Use upquote if available, for straight quotes in verbatim environments
\IfFileExists{upquote.sty}{\usepackage{upquote}}{}
\IfFileExists{microtype.sty}{% use microtype if available
  \usepackage[]{microtype}
  \UseMicrotypeSet[protrusion]{basicmath} % disable protrusion for tt fonts
}{}
\makeatletter
\@ifundefined{KOMAClassName}{% if non-KOMA class
  \IfFileExists{parskip.sty}{%
    \usepackage{parskip}
  }{% else
    \setlength{\parindent}{0pt}
    \setlength{\parskip}{6pt plus 2pt minus 1pt}}
}{% if KOMA class
  \KOMAoptions{parskip=half}}
\makeatother
\usepackage{longtable,booktabs,array}
\usepackage{caption}
\captionsetup[table]{skip=6pt}
\newcounter{none} % for unnumbered tables
\usepackage{calc} % for calculating minipage widths
% Correct order of tables after \paragraph or \subparagraph
\usepackage{etoolbox}
\makeatletter
\patchcmd\longtable{\par}{\if@noskipsec\mbox{}\fi\par}{}{}
\makeatother
% Allow footnotes in longtable head/foot
\IfFileExists{footnotehyper.sty}{\usepackage{footnotehyper}}{\usepackage{footnote}}
\makesavenoteenv{longtable}
\usepackage{graphicx}
\makeatletter
\newsavebox\pandoc@box
\newcommand*\pandocbounded[1]{% scales image to fit in text height/width
  \sbox\pandoc@box{#1}%
  \Gscale@div\@tempa{\textheight}{\dimexpr\ht\pandoc@box+\dp\pandoc@box\relax}%
  \Gscale@div\@tempb{\linewidth}{\wd\pandoc@box}%
  \ifdim\@tempb\p@<\@tempa\p@\let\@tempa\@tempb\fi% select the smaller of both
  \ifdim\@tempa\p@<\p@\scalebox{\@tempa}{\usebox\pandoc@box}%
  \else\usebox{\pandoc@box}%
  \fi%
}
% Set default figure placement to htbp
\def\fps@figure{htbp}
\makeatother
\setlength{\emergencystretch}{3em} % prevent overfull lines
\providecommand{\tightlist}{%
  \setlength{\itemsep}{0pt}\setlength{\parskip}{0pt}}
\usepackage{bookmark}
\IfFileExists{xurl.sty}{\usepackage{xurl}}{} % add URL line breaks if available
\urlstyle{same}
% fallback for those not using the hyperref driver hyperxmp:
\makeatletter
\@ifundefined{xmpquote}{\newcommand{\xmpquote}[1]{#1}}{}
\makeatother
\hypersetup{
  hidelinks,
  pdfcreator={LaTeX via pandoc}}

\author{}
\date{}

\begin{document}

\section{Ruang Hidup yang Terampas}\label{ruang-hidup-yang-terampas}

Analisis dinamika konflik sosial dan alokasi ruang agraria dalam konteks
pembangunan kawasan.

\begin{center}\rule{0.5\linewidth}{0.5pt}\end{center}

Ekspansi industri ekstraktif dan proyek strategis berimplikasi pada
dinamika sosial dan penggunaan lahan masyarakat. Data empiris mencatat
akumulasi \textbf{95 kasus konflik agraria}. Konflik ini berkaitan erat
dengan perubahan tata guna lahan dan alokasi ruang di berbagai daerah.

Aktor dan sektor pemicu konflik mencakup sektor \textbf{Kehutanan}
(Hutan Lindung, Produksi, Konservasi), \textbf{Infrastruktur \& PSN}
(Bendungan, Transmigrasi, Kawasan Industri), hingga proyek
\textbf{Pariwisata \& Pesisir}. Tiga sektor utama (Perkebunan,
Kehutanan, dan Pertambangan) menyumbang porsi \textbf{82.1\%} dari
keseluruhan catatan konflik.

\begin{center}\rule{0.5\linewidth}{0.5pt}\end{center}

\subsubsection{Ringkasan Metrik Agregat}\label{ringkasan-metrik-agregat}

{\def\LTcaptype{none} % do not increment counter
\begin{longtable}[]{@{}
  >{\raggedright\arraybackslash}p{(\linewidth - 4\tabcolsep) * \real{0.3333}}
  >{\raggedright\arraybackslash}p{(\linewidth - 4\tabcolsep) * \real{0.3333}}
  >{\raggedright\arraybackslash}p{(\linewidth - 4\tabcolsep) * \real{0.3333}}@{}}
\toprule\noalign{}
\begin{minipage}[b]{\linewidth}\raggedright
Indikator
\end{minipage} & \begin{minipage}[b]{\linewidth}\raggedright
Nilai
\end{minipage} & \begin{minipage}[b]{\linewidth}\raggedright
Keterangan
\end{minipage} \\
\midrule\noalign{}
\endhead
\bottomrule\noalign{}
\endlastfoot
\textbf{Total Kasus Konflik} & \textbf{95 kasus} & Catatan insiden
sengketa agraria dan tata guna lahan. \\
\textbf{Korban Terdampak (Jiwa)} & \textbf{90,582 jiwa} & Estimasi
jumlah warga yang terdampak oleh konflik sengketa lahan. \\
\textbf{Status: Belum Ditangani} & \textbf{51 kasus} & Kasus sengketa
yang masih dalam proses penanganan. \\
\textbf{Komunitas Terdampak} & \textbf{59 komunitas} & Kelompok tani dan
komunitas lokal yang terlibat dalam sengketa. \\
\textbf{Sektor Perkebunan} & \textbf{25 kasus} & Sengketa tumpang tindih
Hak Guna Usaha (HGU) perkebunan dengan lahan masyarakat. \\
\textbf{Sektor Kehutanan} & \textbf{30 kasus} & Sengketa batas kawasan
hutan produksi dan konservasi dengan wilayah kelola lokal. \\
\textbf{Sektor Pertambangan} & \textbf{23 kasus} & Sengketa alokasi
lahan untuk operasi pertambangan dan fasilitas hilirisasi. \\
\textbf{Infrastruktur \& PSN} & \textbf{13 kasus} & Sengketa pengadaan
tanah untuk Proyek Strategis Nasional. \\
\textbf{Pariwisata \& Pesisir} & \textbf{3 kasus} & Sengketa pemanfaatan
wilayah pesisir dan kawasan pariwisata. \\
\textbf{Keterlibatan Pemerintah} & \textbf{69 kasus} & Keterlibatan
instansi pemerintah dalam fasilitasi atau sengketa lahan. \\
\textbf{Keterlibatan Korporasi} & \textbf{52 kasus} & Entitas BUMN atau
swasta yang terlibat dalam sengketa lahan. \\
\end{longtable}
}

\emph{Sumber Analisis Data: Konsorsium Pembaruan Agraria (KPA) / Tanah
Kita}

\begin{center}\rule{0.5\linewidth}{0.5pt}\end{center}

\subsubsection{4.1 Tren Eskalasi Konflik Agraria Seiring Ekspansi
Industri}\label{tren-eskalasi-konflik-agraria-seiring-ekspansi-industri}

Visualisasi \emph{time-series} di bawah ini memberikan gambaran korelasi
antara ekspansi industri dan dinamika konflik agraria di daratan
Sulawesi. Secara historis, perbandingan dua periode waktu menunjukkan
perbedaan tingkat insidensi konflik. Pada periode pra-2005, sistem
pendataan mencatat \textbf{13 kasus} konflik agraria.

Pada periode pasca-2005 hingga saat ini, data mencatat \textbf{82 kasus}
konflik lahan, yang mencerminkan peningkatan sebesar \textbf{630.8\%}
dibandingkan periode sebelumnya. Perubahan tren ini beriringan dengan
penerbitan Izin Usaha Pertambangan (IUP) serta ekspansi Hak Guna Usaha
(HGU) untuk perkebunan.

Penelusuran tren satu dekade terakhir menunjukkan bahwa sengketa agraria
mencakup berbagai sektor, termasuk pertambangan nikel, infrastruktur,
dan Proyek Strategis Nasional. Akumulasi \textbf{95 insiden historis}
ini mengindikasikan perlunya tata kelola alokasi lahan dan perlindungan
hak masyarakat lokal yang lebih seimbang di kawasan investasi.

\begin{figure}
\centering
\pandocbounded{\includegraphics[keepaspectratio,alt={Peningkatan Signifikan Konflik Agraria di Sulawesi (1990-2025)}]{visuals_bab4/chart_4_1_konflik_timeseries.png}}
\caption{Peningkatan Signifikan Konflik Agraria di Sulawesi (1990-2025)}
\end{figure}

\begin{quote}
\textbf{Interpretasi Ekologis: Puncak Insidensi Konflik 2017}

Grafik memperlihatkan peningkatan insidensi konflik yang memuncak pada
\textbf{tahun 2017} dengan \textbf{75 kasus konflik}. Pembedahan data
sektoral menunjukkan konsentrasi pada sektor \textbf{Kehutanan (40
kasus)} dan \textbf{Perkebunan (21 kasus)}, diikuti oleh
\textbf{Pertambangan dan Infrastruktur PSN}. Periode ini bertepatan
dengan percepatan pelepasan kawasan hutan dan Izin Pinjam Pakai Kawasan
Hutan (IPPKH) untuk mendukung proyek strategis dan kawasan industri.
\end{quote}

\begin{quote}
\textbf{Interpretasi Ekologis dan Sosial:}

Peningkatan insidensi konflik beririsan dengan dinamika perizinan
kawasan. Pengelolaan alokasi ruang dan perlindungan hak masyarakat di
wilayah investasi menjadi faktor penting untuk meminimalkan dampak
sosial.
\end{quote}

\begin{center}\rule{0.5\linewidth}{0.5pt}\end{center}

\subsubsection{4.2 Sebaran Sektoral: Dampak Masyarakat dan Penggunaan
Lahan}\label{sebaran-sektoral-dampak-masyarakat-dan-penggunaan-lahan}

Visualisasi komparatif di bawah ini menggambarkan skala dampak sosial
dan penggunaan lahan berdasarkan sektor industri. Data menunjukkan bahwa
\textbf{Sektor Kehutanan} mencatatkan jumlah warga terdampak sebanyak
\textbf{21,886 jiwa}, berkaitan dengan tumpang tindih kawasan hutan
produksi, konservasi, dan Hutan Tanaman Industri (HTI) dengan wilayah
kelola masyarakat lokal.

Menyusul berikutnya adalah \textbf{Sektor Pertambangan} dengan total
korban terdampak sebanyak \textbf{54,658 jiwa}, yang beririsan dengan
proyek hilirisasi nikel dan tambang terbuka di kawasan pesisir dan
pertanian.

Dari dimensi penggunaan lahan (luasan hektar yang terlibat sengketa),
\textbf{Sektor Perkebunan} mencatatkan luas sengketa terbesar yaitu
\textbf{77,902 Hektar}, disusul oleh sektor Kehutanan seluas
\textbf{66,193 Ha} dan Pertambangan seluas \textbf{441,286 Ha}. Data ini
menunjukkan bahwa dinamika penguasaan lahan di tiga sektor tersebut
berkorelasi dengan tingginya insidensi sengketa agraria di tingkat
lokal.

\begin{figure}
\centering
\pandocbounded{\includegraphics[keepaspectratio,alt={Peningkatan Signifikan Korban Terdampak (Jiwa) per Tahun}]{visuals_bab4/chart_4_2a_jiwa.png}}
\caption{Peningkatan Signifikan Korban Terdampak (Jiwa) per Tahun}
\end{figure}

\begin{figure}
\centering
\pandocbounded{\includegraphics[keepaspectratio,alt={Monopoli Area Konflik (Hektar) per Tahun}]{visuals_bab4/chart_4_2b_ha.png}}
\caption{Monopoli Area Konflik (Hektar) per Tahun}
\end{figure}

\begin{quote}
\textbf{Interpretasi Ekologis dan Sosial:}

Dinamika Grafik mencerminkan akumulasi dampak sosial di wilayah industri
yang memerlukan perhatian dalam pengelolaan sengketa lahan.
\end{quote}

\subsubsection{Bedah Forensik Anomali (Spike) Konflik
Agraria}\label{bedah-forensik-anomali-spike-konflik-agraria}

Berdasarkan ekstraksi dataset secara mendalam, berikut adalah bedah
anatomis dari lonjakan-lonjakan ekstrem (\emph{spikes}) yang terjadi
pada grafik \textbf{Peningkatan Signifikan Korban Terdampak (Jiwa)} dan
\textbf{Monopoli Area Konflik (Hektar)} di wilayah ini.

\paragraph{ANOMALI JIWA 1: Lonjakan Korban Jiwa Tahun
2022}\label{anomali-jiwa-1-lonjakan-korban-jiwa-tahun-2022}

\begin{itemize}
\tightlist
\item
  \textbf{Kasus Utama Pendongkrak Statistik:} Konflik Koalisi Selamatkan
  Pulau Wawoni
\item
  \textbf{Total Korban (Kasus Ini):} 37,000 Jiwa
\item
  \textbf{Perusahaan Terlibat:} PT Gema Kreasi Perdana (GKP)
\item
  \textbf{Narasi Singkat:} Keberadaan PT Gema Kreasi Perdana (GKP), anak
  perusahaan Harita Group, di Pulau Wawonii, Kabupaten Konawe Kepulauan,
  Sulawesi Tenggara mengancam keselamatan warga, menggusur lahan
  pertanian-perkebunan dan berpotensi mencemari laut. Keduanya merupakan
  sumber ekonomi utama dari lebih dari 37.000 jiwa warga di Pulau
  Wawonii. Tidak hanya mengancam keselamatan dan ruang hidup warga,
  aktivitas pertambangan yang dilakukan PT GKP adalah ilegal karena
  ban\ldots{}
\item
  \textbf{Sumber Referensi:} Laporan WALHI
  (\href{https://www.google.com/search?q=Konflik\%20Koalisi\%20Selamatkan\%20Pulau\%20Wawoni\%20PT\%20Gema\%20Kreasi\%20Perdana\%20\%28GKP\%29}{Telusuri
  Berita Kasus} \textbar{}
  \href{https://tanahkita.id/data/konflik/detil/SWpxcGxoUThtUVk}{Link
  Asli TanahKita})
\end{itemize}

\paragraph{ANOMALI JIWA 2: Lonjakan Korban Jiwa Tahun
2012}\label{anomali-jiwa-2-lonjakan-korban-jiwa-tahun-2012}

\begin{itemize}
\tightlist
\item
  \textbf{Kasus Utama Pendongkrak Statistik:} Konflik Masyarakat Tumpang
  Pitu dengan PT BSI
\item
  \textbf{Total Korban (Kasus Ini):} 13,936 Jiwa
\item
  \textbf{Perusahaan Terlibat:} PT Bumi Suksindo \textbar{} PT. Merdeka
  Copper Gold
\item
  \textbf{Narasi Singkat:} Gunung Tumpang Pitu yang masuk wilayah Desa
  Sumberagung, Kecamatan Pesanggaran, Kabupaten Banyuwangi, Jawa Timur,
  merupakan lokasi tambang emas dari PT Bumi Suksindo ( PT BSI) seluas
  4.998 Ha, melalui Izin Usaha Pertambangan Operasi Produksi (IUP OP)
  berdasarkan Keputusan Bupati Banyuwangi No.~188/547/KEP/429.011/2012
  tanggal 9 Juli 2012 sebagaimana terakhir kali diubah dengan Keputusan
  Bupati Banyuwangi, No.~188/928/KEP/429.011/2012 tertanggal 7\ldots{}
\item
  \textbf{Sumber Referensi:} Laporan WALHI JATIM, Front Nahdliyin Untuk
  Kedaulatan Sumber Daya Alam
  (\href{https://www.google.com/search?q=Konflik\%20Masyarakat\%20Tumpang\%20Pitu\%20dengan\%20PT\%20BSI\%20PT\%20Bumi\%20Suksindo\%20\%7C\%20PT.\%20Merdeka\%20Copper\%20Gold}{Telusuri
  Berita Kasus} \textbar{}
  \href{https://tanahkita.id/data/konflik/detil/WWg2QmItRXZ4SGs}{Link
  Asli TanahKita})
\end{itemize}

\paragraph{ANOMALI AREA 1: Monopoli Area Konflik Tahun
2017}\label{anomali-area-1-monopoli-area-konflik-tahun-2017}

\begin{itemize}
\tightlist
\item
  \textbf{Kasus Utama Pendongkrak Statistik:} Konflik Nelayan Desa
  Boddia dengan Pertambangan Pasir Laut PT. Boskalis dan PT. Jan De Nul
\item
  \textbf{Total Daratan Dirampas (Kasus Ini):} 146,201 Hektar
\item
  \textbf{Perusahaan Terlibat:} PT. Boskalis \textbar{} PT. Jan De Nul
\item
  \textbf{Narasi Singkat:} Pertambangan pasir laut yang dilakukan oleh
  PT. Boskalis dan PT. Jan De Null menjadi dasar dari munculnya konflik
  yang ada di Desa Boddia, Khususnya di Dusun Boddia dan Manjalling.
  Pertambangan pasir laut ini merupakan bagian dari proyek reklamasi
  yang telah disepakati oleh Gubernur dan DPRD Provinsi Sulawesi Selatan
  di dalam Ranperda RZWP3K Sulsel. Kebutuhan reklamasi ini menciptakan
  beberapa daerah yang dialokasikan menjadi zona tambang pasir l\ldots{}
\item
  \textbf{Sumber Referensi:} Laporan WALHI Sulawesi Selatan
  (\href{https://www.google.com/search?q=Konflik\%20Nelayan\%20Desa\%20Boddia\%20dengan\%20Pertambangan\%20Pasir\%20Laut\%20PT.\%20Boskalis\%20dan\%20PT.\%20Jan\%20De\%20Nul\%20PT.\%20Boskalis\%20\%7C\%20PT.\%20Jan\%20De\%20Nul}{Telusuri
  Berita Kasus} \textbar{}
  \href{https://tanahkita.id/data/konflik/detil/dXpRbGVnRTJHS2c}{Link
  Asli TanahKita})
\end{itemize}

\paragraph{ANOMALI AREA 2: Monopoli Area Konflik Tahun
2022}\label{anomali-area-2-monopoli-area-konflik-tahun-2022}

\begin{itemize}
\tightlist
\item
  \textbf{Kasus Utama Pendongkrak Statistik:} Konflik Koalisi Selamatkan
  Pulau Wawoni
\item
  \textbf{Total Daratan Dirampas (Kasus Ini):} 86,758 Hektar
\item
  \textbf{Perusahaan Terlibat:} PT Gema Kreasi Perdana (GKP)
\item
  \textbf{Narasi Singkat:} Keberadaan PT Gema Kreasi Perdana (GKP), anak
  perusahaan Harita Group, di Pulau Wawonii, Kabupaten Konawe Kepulauan,
  Sulawesi Tenggara mengancam keselamatan warga, menggusur lahan
  pertanian-perkebunan dan berpotensi mencemari laut. Keduanya merupakan
  sumber ekonomi utama dari lebih dari 37.000 jiwa warga di Pulau
  Wawonii. Tidak hanya mengancam keselamatan dan ruang hidup warga,
  aktivitas pertambangan yang dilakukan PT GKP adalah ilegal karena
  ban\ldots{}
\item
  \textbf{Sumber Referensi:} Laporan WALHI
  (\href{https://www.google.com/search?q=Konflik\%20Koalisi\%20Selamatkan\%20Pulau\%20Wawoni\%20PT\%20Gema\%20Kreasi\%20Perdana\%20\%28GKP\%29}{Telusuri
  Berita Kasus} \textbar{}
  \href{https://tanahkita.id/data/konflik/detil/SWpxcGxoUThtUVk}{Link
  Asli TanahKita})
\end{itemize}

\begin{center}\rule{0.5\linewidth}{0.5pt}\end{center}

\subsubsection{4.3 Indikasi Represi dan Kriminalisasi dalam Konflik
Agraria}\label{indikasi-represi-dan-kriminalisasi-dalam-konflik-agraria}

Data kuantitatif di wilayah Sulawesi mencatat indikasi terjadinya
represi dan tindakan kriminalisasi dalam sebagian sengketa agraria. Dari
database yang didokumentasikan, terdapat \textbf{46 kasus indikasi
kriminalisasi} dan \textbf{93 warga/aktivis lingkungan yang tercatat
pernah ditangkap} dalam penanganan sengketa lahan.

Berdasarkan distribusi sektoral, \textbf{Sektor Pertambangan}
mencatatkan frekuensi indikasi represi tertinggi dengan \textbf{17
kasus}. Tahun dengan jumlah catatan insiden represi tertinggi adalah
\textbf{2017} dengan \textbf{7 kasus}.

Catatan ini menunjukkan pentingnya pendekatan hukum yang adil,
penyelesaian konflik secara ramah HAM, serta jaminan perlindungan bagi
pejuang lingkungan dan komunitas lokal sesuai dengan peraturan
perundang-undangan.

{\def\LTcaptype{none} % do not increment counter
\begin{longtable}[]{@{}
  >{\raggedright\arraybackslash}p{(\linewidth - 6\tabcolsep) * \real{0.2500}}
  >{\raggedright\arraybackslash}p{(\linewidth - 6\tabcolsep) * \real{0.2500}}
  >{\raggedright\arraybackslash}p{(\linewidth - 6\tabcolsep) * \real{0.2500}}
  >{\raggedright\arraybackslash}p{(\linewidth - 6\tabcolsep) * \real{0.2500}}@{}}
\toprule\noalign{}
\begin{minipage}[b]{\linewidth}\raggedright
Kasus Indikasi Kriminalisasi
\end{minipage} & \begin{minipage}[b]{\linewidth}\raggedright
Warga/Aktivis Ditangkap
\end{minipage} & \begin{minipage}[b]{\linewidth}\raggedright
Korban Luka-luka
\end{minipage} & \begin{minipage}[b]{\linewidth}\raggedright
Korban Tewas
\end{minipage} \\
\midrule\noalign{}
\endhead
\bottomrule\noalign{}
\endlastfoot
\textbf{46 Kasus} & \textbf{93 Orang} & \textbf{4 Orang} & \textbf{1
Orang} \\
\end{longtable}
}

\begin{figure}
\centering
\pandocbounded{\includegraphics[keepaspectratio,alt={Tren Kasus Kriminalisasi \& Represi (Pasca 2000)}]{visuals_bab4/chart_4_3a_kriminalisasi_trend.png}}
\caption{Tren Kasus Kriminalisasi \& Represi (Pasca 2000)}
\end{figure}

\begin{figure}
\centering
\pandocbounded{\includegraphics[keepaspectratio,alt={Sektor Industri Paling Represif}]{visuals_bab4/chart_4_3b_sektor_represif.png}}
\caption{Sektor Industri Paling Represif}
\end{figure}

\begin{quote}
\textbf{Interpretasi Ekologis \& Hak Asasi Manusia:} Keberadaan kasus
kriminalisasi di sekitar area konsesi (terutama Pertambangan)
mengindikasikan pentingnya jaminan perlindungan ruang sipil dan
penghormatan HAM dalam setiap proses pembangunan.
\end{quote}

\begin{center}\rule{0.5\linewidth}{0.5pt}\end{center}

\subsubsection{4.4 Pembuktian Statistik: Ekspansi vs Eskalasi
Konflik}\label{pembuktian-statistik-ekspansi-vs-eskalasi-konflik}

Hipotesis utama dalam evaluasi ini adalah bahwa \textbf{industrialisasi
dan ekspansi korporasi} berbanding lurus dengan \textbf{eskalasi konflik
dan represi} terhadap masyarakat. Untuk mengujinya secara statistik
sesuai pedoman D3TLH, analisis dibagi menjadi dua bagian: (1) Komparasi
metrik Before-After, dan (2) Uji signifikansi Crosstab Chi-Square. Unit
observasinya adalah catatan kejadian letupan konflik historis.

\paragraph{A. Analisis Komparatif Before-After (Pra vs Era
Hilirisasi)}\label{a.-analisis-komparatif-before-after-pra-vs-era-hilirisasi}

Perbandingan absolut eskalasi konflik agraria sebelum dan sesudah rezim
hilirisasi masif dimulai (cut-off tahun 2014).

{\def\LTcaptype{none} % do not increment counter
\begin{longtable}[]{@{}
  >{\raggedright\arraybackslash}p{(\linewidth - 8\tabcolsep) * \real{0.2000}}
  >{\raggedright\arraybackslash}p{(\linewidth - 8\tabcolsep) * \real{0.2000}}
  >{\raggedright\arraybackslash}p{(\linewidth - 8\tabcolsep) * \real{0.2000}}
  >{\raggedright\arraybackslash}p{(\linewidth - 8\tabcolsep) * \real{0.2000}}
  >{\raggedright\arraybackslash}p{(\linewidth - 8\tabcolsep) * \real{0.2000}}@{}}
\toprule\noalign{}
\begin{minipage}[b]{\linewidth}\raggedright
Periode
\end{minipage} & \begin{minipage}[b]{\linewidth}\raggedright
Rata-rata Konflik
\end{minipage} & \begin{minipage}[b]{\linewidth}\raggedright
Total Letupan
\end{minipage} & \begin{minipage}[b]{\linewidth}\raggedright
Warga Ditangkap
\end{minipage} & \begin{minipage}[b]{\linewidth}\raggedright
Korban Tewas
\end{minipage} \\
\midrule\noalign{}
\endhead
\bottomrule\noalign{}
\endlastfoot
\textbf{Pra-Ekspansi (1990 - 2013)} & \textbf{2.1 Kasus/Tahun} & 37
kejadian & 83 jiwa & 0 jiwa \\
\textbf{Pasca-Ekspansi (2014 - 2024)} & \textbf{5.5 Kasus/Tahun} & 55
kejadian & 10 jiwa & 1 jiwa \\
\end{longtable}
}

\paragraph{B. Uji Statistik Crosstab
(Chi-Square)}\label{b.-uji-statistik-crosstab-chi-square}

\textbf{Variabel Independen (X):} Periode Ekspansi Industri

\textbf{Variabel Dependen (Y):} Tingkat Represi \& Kriminalisasi

\paragraph{Case Processing Summary}\label{case-processing-summary}

{\def\LTcaptype{none} % do not increment counter
\begin{longtable}[]{@{}
  >{\raggedright\arraybackslash}p{(\linewidth - 12\tabcolsep) * \real{0.1429}}
  >{\raggedright\arraybackslash}p{(\linewidth - 12\tabcolsep) * \real{0.1429}}
  >{\raggedright\arraybackslash}p{(\linewidth - 12\tabcolsep) * \real{0.1429}}
  >{\raggedright\arraybackslash}p{(\linewidth - 12\tabcolsep) * \real{0.1429}}
  >{\raggedright\arraybackslash}p{(\linewidth - 12\tabcolsep) * \real{0.1429}}
  >{\raggedright\arraybackslash}p{(\linewidth - 12\tabcolsep) * \real{0.1429}}
  >{\raggedright\arraybackslash}p{(\linewidth - 12\tabcolsep) * \real{0.1429}}@{}}
\toprule\noalign{}
\begin{minipage}[b]{\linewidth}\raggedright
\end{minipage} & \begin{minipage}[b]{\linewidth}\raggedright
Valid N
\end{minipage} & \begin{minipage}[b]{\linewidth}\raggedright
Valid \%
\end{minipage} & \begin{minipage}[b]{\linewidth}\raggedright
Missing N
\end{minipage} & \begin{minipage}[b]{\linewidth}\raggedright
Missing \%
\end{minipage} & \begin{minipage}[b]{\linewidth}\raggedright
Total N
\end{minipage} & \begin{minipage}[b]{\linewidth}\raggedright
Total \%
\end{minipage} \\
\midrule\noalign{}
\endhead
\bottomrule\noalign{}
\endlastfoot
Periode Ekspansi Industri * Tingkat Represi \& Kriminalisasi & 523 &
100.0\% & 0 & 0.0\% & 523 & 100.0\% \\
\end{longtable}
}

\paragraph{Periode Ekspansi Industri * Tingkat Represi \& Kriminalisasi
Crosstabulation}\label{periode-ekspansi-industri-tingkat-represi-kriminalisasi-crosstabulation}

{\def\LTcaptype{none} % do not increment counter
\begin{longtable}[]{@{}llll@{}}
\toprule\noalign{}
& Baseline (Tanpa Kriminalisasi) & Ada Represi/Kriminalisasi & Total \\
\midrule\noalign{}
\endhead
\bottomrule\noalign{}
\endlastfoot
\textbf{Pra-ekspansi (\textless{} 2014)} Count & 179 & 73 & 252 \\
\textbf{Pra-ekspansi (\textless{} 2014)} Expected & 157.6 & 94.4 &
252.0 \\
\textbf{Pasca-ekspansi (≥ 2014)} Count & 148 & 123 & 271 \\
\textbf{Pasca-ekspansi (≥ 2014)} Expected & 169.4 & 101.6 & 271.0 \\
\textbf{Total} Count & 327 & 196 & 523 \\
\textbf{Total} Expected & 327.0 & 196.0 & 523.0 \\
\end{longtable}
}

\paragraph{Chi-Square Tests}\label{chi-square-tests}

\textbf{Periode Ekspansi Industri * Tingkat Represi \& Kriminalisasi}

{\def\LTcaptype{none} % do not increment counter
\begin{longtable}[]{@{}llll@{}}
\toprule\noalign{}
& Value & df & Asymp. Sig. (2-sided) \\
\midrule\noalign{}
\endhead
\bottomrule\noalign{}
\endlastfoot
Pearson Chi-Square & 14.331 & 1 & 0.000 \\
Likelihood Ratio & 14.447 & 1 & 0.000 \\
Linear-by-Linear Association & 14.995 & 1 & 0.000 \\
N of Valid Cases & 523 & & \\
\end{longtable}
}

\subsubsection{Ringkasan Uji Hipotesis}\label{ringkasan-uji-hipotesis}

\textbf{Result: SIGNIFIKAN}

{\def\LTcaptype{none} % do not increment counter
\begin{longtable}[]{@{}ll@{}}
\toprule\noalign{}
Parameter & Nilai \\
\midrule\noalign{}
\endhead
\bottomrule\noalign{}
\endlastfoot
P-Value & 0.0002 \\
Chi-Square & 14.331 \\
df & 1 \\
\textbf{Odds Ratio (Risk Estimate)} & \textbf{2.038} \\
\end{longtable}
}

\begin{quote}
\textbf{Interpretasi Sosial Kritis:} Temuan ini sangat krusial:
pergeseran status \textbf{Periode Ekspansi Industri} terbukti
\textbf{berkorelasi kuat dan signifikan} dengan \textbf{Tingkat Represi
\& Kriminalisasi} (P \textless{} 0.05). Angka Odds Ratio (OR: 2.038)
menjadi konfirmasi empiris bahwa narasi hilirisasi dan investasi
bukanlah agenda nirkekerasan---ekspansi spasial mereka mutlak
mengeskalasi pelanggaran hak asasi masyarakat tapak.
\end{quote}

\begin{center}\rule{0.5\linewidth}{0.5pt}\end{center}

\subsubsection{Ringkasan Eksekutif Seluruh Skenario
Crosstab}\label{ringkasan-eksekutif-seluruh-skenario-crosstab}

Tabel di bawah ini merangkum hasil pengujian statistik (Chi-Square)
untuk semua kemungkinan kombinasi antara indikator Ekspansi (X) dan
Eskalasi Konflik (Y) pada panel data yang sama.

{\def\LTcaptype{none} % do not increment counter
\begin{longtable}[]{@{}
  >{\raggedright\arraybackslash}p{(\linewidth - 10\tabcolsep) * \real{0.1667}}
  >{\raggedright\arraybackslash}p{(\linewidth - 10\tabcolsep) * \real{0.1667}}
  >{\raggedright\arraybackslash}p{(\linewidth - 10\tabcolsep) * \real{0.1667}}
  >{\raggedright\arraybackslash}p{(\linewidth - 10\tabcolsep) * \real{0.1667}}
  >{\raggedright\arraybackslash}p{(\linewidth - 10\tabcolsep) * \real{0.1667}}
  >{\raggedright\arraybackslash}p{(\linewidth - 10\tabcolsep) * \real{0.1667}}@{}}
\toprule\noalign{}
\begin{minipage}[b]{\linewidth}\raggedright
Variabel Independen (X)
\end{minipage} & \begin{minipage}[b]{\linewidth}\raggedright
Variabel Dependen (Y)
\end{minipage} & \begin{minipage}[b]{\linewidth}\raggedright
Chi-Square
\end{minipage} & \begin{minipage}[b]{\linewidth}\raggedright
P-Value
\end{minipage} & \begin{minipage}[b]{\linewidth}\raggedright
Odds Ratio
\end{minipage} & \begin{minipage}[b]{\linewidth}\raggedright
Kesimpulan
\end{minipage} \\
\midrule\noalign{}
\endhead
\bottomrule\noalign{}
\endlastfoot
Periode Ekspansi Industri & Tingkat Represi \& Kriminalisasi & 14.331 &
0.000 & 2.04 & 🟢 SIGNIFIKAN \\
Periode Ekspansi Industri & Tingkat Penelantaran Kasus & 3.283 & 0.070 &
1.40 & 🔴 TIDAK SIGNIFIKAN \\
Periode Ekspansi Industri & Tingkat Insiden Fisik (Luka/Tewas/Ditangkap)
& 2.636 & 0.104 & 0.53 & 🔴 TIDAK SIGNIFIKAN \\
Tipe Sektor (Tambang vs Non-Tambang) & Tingkat Represi \& Kriminalisasi
& 13.609 & 0.000 & 2.84 & 🟢 SIGNIFIKAN \\
Tipe Sektor (Tambang vs Non-Tambang) & Tingkat Penelantaran Kasus &
4.742 & 0.029 & 1.89 & 🟢 SIGNIFIKAN \\
Tipe Sektor (Tambang vs Non-Tambang) & Tingkat Insiden Fisik
(Luka/Tewas/Ditangkap) & 0.665 & 0.415 & 1.66 & 🔴 TIDAK SIGNIFIKAN \\
Keterlibatan Aparat/Pemerintah & Tingkat Represi \& Kriminalisasi &
55.633 & 0.000 & 4.80 & 🟢 SIGNIFIKAN \\
Keterlibatan Aparat/Pemerintah & Tingkat Penelantaran Kasus & 1.135 &
0.287 & 0.81 & 🔴 TIDAK SIGNIFIKAN \\
Keterlibatan Aparat/Pemerintah & Tingkat Insiden Fisik
(Luka/Tewas/Ditangkap) & 3.253 & 0.071 & 2.23 & 🔴 TIDAK SIGNIFIKAN \\
\end{longtable}
}

\begin{quote}
\textbf{Pembedahan Realitas Kemanusiaan:}

Dari \textbf{9 skenario pengujian}, terdapat \textbf{4 skenario yang
terbukti SIGNIFIKAN}.
\end{quote}

Angka-angka pada tabel di atas bukan sekadar statistik di atas kertas,
melainkan \textbf{bukti empiris} dari brutalitas pembangunan. Tingginya
angka kemunculan represi pada skenario yang signifikan menegaskan bahwa
setiap kali wilayah operasi investasi diperlebar, probabilitas
dihadapkannya moncong senjata kepada warga melonjak drastis.

Skenario yang \emph{TIDAK SIGNIFIKAN} tidak berarti rezim terbebas dari
dosa kekerasan, melainkan bukti bahwa represi terhadap warga yang
mempertahankan tanahnya telah menjadi kultur mapan yang menyebar secara
sporadis melampaui sekat waktu dan korporasi.

\begin{center}\rule{0.5\linewidth}{0.5pt}\end{center}

\subsubsection{4.5 Peta Entitas Aktor: Korporasi dan Organisasi
Masyarakat}\label{peta-entitas-aktor-korporasi-dan-organisasi-masyarakat}

Analisis entitas aktor berbasis pemrosesan teks (\emph{string parsing})
terhadap catatan kronologi dokumentasi TanahKita memetakan keterlibatan
berbagai pihak dalam sengketa agraria. Hasil ekstraksi teks
mengidentifikasi entitas korporasi, lembaga pemerintah, serta organisasi
masyarakat sipil yang tercatat dalam dokumentasi kasus. Grafik frekuensi
di bawah menampilkan entitas korporasi dan kelompok masyarakat yang
paling sering teridentifikasi dalam catatan sengketa lahan.

\paragraph{Top 10 Entitas Korporasi Paling
Dominan}\label{top-10-entitas-korporasi-paling-dominan}

\begin{figure}
\centering
\pandocbounded{\includegraphics[keepaspectratio,alt={Top 10 Entitas Korporasi Paling Dominan}]{visuals_bab4/chart_4_5a_korporasi.png}}
\caption{Top 10 Entitas Korporasi Paling Dominan}
\end{figure}

\begin{quote}
\textbf{Analisis Data Korporasi:} Ekstraksi teks mencatat frekuensi
penyebutan entitas \textbf{PT Perkebunan Nusantara (PTPN)} dalam
\textbf{153 catatan kasus terpisah}.
\end{quote}

\paragraph{Top Aktor Proksi \& Vigilante
Terdeteksi}\label{top-aktor-proksi-vigilante-terdeteksi}

\begin{figure}
\centering
\pandocbounded{\includegraphics[keepaspectratio,alt={Top Aktor Proksi \& Vigilante Terdeteksi}]{visuals_bab4/chart_4_5b_vigilante.png}}
\caption{Top Aktor Proksi \& Vigilante Terdeteksi}
\end{figure}

\begin{quote}
\textbf{Analisis Kritis Proksi/Vigilante:} Kemunculan kelompok sipil
seperti \textbf{Preman} (terdeteksi hingga \textbf{16 kali}) menangkap
besarnya skala orkestrasi horizontal. Korporasi seringkali menggunakan
jasa pengamanan swakarsa, kelompok preman, hingga ormas vigilante
sebagai ``bemper proksi'' untuk mengintimidasi warga lokal dan memecah
belah solidaritas akar rumput.
\end{quote}

\emph{* Grafik di atas hanya menampilkan Top 10 entitas. Untuk melihat
daftar lengkap dan detail seluruh aktor yang terdeteksi, silakan buka
tabel data.}

\end{document}
